% Chapter 8: Observability, Drift Monitoring & Active Learning

\begin{frame}[plain]
    \vfill
    \centering
    {\footnotesize \textbf{\textcolor{deepblue}{\textsc{TaxonVision-MLOps Architecture}}}} \\[0.35cm]
    {\large \textbf{\textcolor{signalorange}{CHAPTER 08}}} \\[0.25cm]
    {\huge \textbf{Observability \& Active Learning}} \\[0.35cm]
    {\normalsize \textcolor{black!70}{Prometheus Telemetry, Concept Drift Detection \& Continuous Human-in-the-Loop Triage}} \\[0.5cm]
    \textcolor{deepblue}{\rule{0.55\linewidth}{1.2pt}} \\[0.5cm]
    \statusdone
    \vfill
    \script{In Chapter 8, we present production observability: tracking domain-specific epistemic metrics, detecting embedding drift via Wasserstein distance, and closing the active learning feedback loop.}
    \qa{How does MLOps observability differ from traditional software monitoring?}{An ML API can return HTTP 200 OK while failing silently; we monitor average conformal prediction set sizes and drift distances to detect data shifts in real time.}
\end{frame}

\begin{frame}{Observability}{Production Telemetry: Prometheus Metrics \& Grafana}
    \begin{block}{Domain-Specific Telemetry Beyond HTTP 200}
        \begin{itemize}\small
            \item \textbf{System Metrics:} p95/p99 inference latency ($<50$ ms budget), GPU VRAM utilization, request throughput.
            \item \textbf{Epistemic Metrics:} Average conformal set size ($|C(X)|$), energy-based OOD anomaly rate.
            \item \textbf{Operational Metrics:} Triage referral rate to human taxonomists (Label Studio queue depth).
            \item \textbf{Decoupled Polling:} FastAPI \texttt{/metrics} endpoint scraped every 15s by Prometheus, feeding real-time Grafana dashboards without hot-path inference overhead.
        \end{itemize}
    \end{block}

    \vspace{0.15cm}
    \statusdone

    \script{Observability in MLOps is fundamentally different from traditional DevOps. In web development, you monitor CPU and HTTP status codes. In MLOps, an API can return HTTP 200 OK all day while predicting complete nonsense! We monitor domain-specific metrics: average conformal prediction set size, percentage of samples routed to human review, and OOD energy distributions. If average set size jumps from 1.2 to 4.5, we know immediately that the data distribution has shifted.}
    \qa{How does Prometheus scrape metrics from our service?}{FastAPI instruments latency and counters using the \texttt{prometheus\_client} library, exposing a \texttt{/metrics} endpoint that the Prometheus server polls every 15 seconds.}
    \qa{Why not push metrics directly to a database?}{The pull-based model of Prometheus ensures that if the monitoring system goes down, the inference service is unaffected. It drops the scraping request rather than blocking the main thread on a database timeout.}
\end{frame}

\begin{frame}{Telemetry Architecture}{FastAPI $\rightarrow$ Prometheus $\rightarrow$ Grafana Pipeline}
    \begin{center}
    \resizebox{0.85\textwidth}{!}{
    \begin{tikzpicture}[
        node distance=1.5cm and 2cm,
        box/.style={rectangle, draw=deepblue, thick, fill=blue!5, text width=3cm, align=center, minimum height=1.5cm, rounded corners},
        arrow/.style={->, thick, >=stealth},
        db/.style={cylinder, draw=soilbrown, thick, aspect=0.25, fill=orange!10, text width=2.5cm, align=center, minimum height=1.5cm, shape border rotate=90}
    ]
        % Nodes
        \node[box] (fastapi) {\textbf{FastAPI Pod} \\ \texttt{/metrics} endpoint \\ (Counters, Gauges, Histograms)};
        \node[db, right=of fastapi] (prometheus) {\textbf{Prometheus} \\ Time-Series Database \\ (Scrapes every 15s)};
        \node[box, right=of prometheus] (grafana) {\textbf{Grafana} \\ Real-time Dashboards \\ (Latency, VRAM, Drift)};
        \node[box, below=of fastapi, fill=green!5, draw=plantgreen] (model) {\textbf{ONNX Runtime} \\ Inference Engine};

        % Arrows
        \draw[arrow] (model.west) |- node[near start, left, align=right, xshift=-0.1cm] {Latency \& \\ OOD stats} (fastapi.west);
        \draw[arrow] (prometheus) -- node[above] {HTTP Pull} (fastapi);
        \draw[arrow] (grafana) -- node[above] {PromQL Queries} (prometheus);
    \end{tikzpicture}
    }
    \end{center}

    \script{This diagram illustrates our telemetry flow. The core inference engine passes latency and out-of-distribution statistics up to the FastAPI layer, which aggregates them in memory using the Prometheus client. The Prometheus server periodically pulls these metrics, and Grafana queries Prometheus to build the dashboards you see in the control room.}
    \qa{Why does Prometheus pull rather than the service pushing metrics?}{Pull-based monitoring prevents the application from hanging if the metrics server goes down or becomes slow, ensuring zero impact on the critical inference path.}
\end{frame}

\begin{frame}{Drift Monitoring}{Data \& Concept Drift: Wasserstein-1 Distance vs MMD}
    \begin{block}{Multi-Scale Distributional Drift Detection}
        \begin{itemize}\small
            \item \textbf{Environmental Shift:} Phenological seasonality and new camera sensors cause feature representations to drift from baseline distributions.
            \item \textbf{Online Marginal Tests (1D):} 1D projection with Wasserstein-1 (Earth Mover's) distance for real-time alerting with negligible CPU cost.
            \item \textbf{Offline Manifold Tests (Multivariate):} Maximum Mean Discrepancy (MMD) with RBF kernels evaluates complete structural manifold shifts in scheduled batch audits ($O(N^2)$).
            \item \textbf{Drift Retraining Triggers:} Sustained statistical divergence triggers automated active learning data extraction.
        \end{itemize}
    \end{block}

    \vspace{0.15cm}
    \statusdone

    \script{Machine learning models degrade over time because the real world changes—we call this data drift. If autumn arrives, leaves turn yellow and lighting shifts, but our training set was collected in summer. We use the Wasserstein-1 Earth Mover's Distance to monitor the distribution of incoming feature embeddings against our baseline training distribution. When the statistical distance crosses our alert threshold, the system flags the drift and prepares an active learning cycle.}
    \qa{Why not run multivariate MMD on every single incoming inference request?}{Computing full-manifold MMD across high-dimensional embeddings has quadratic complexity $O(N^2)$ in sample size. We perform 1D Wasserstein distance on projected features in real time and reserve full MMD for periodic offline batch audits.}
\end{frame}

\begin{frame}{Active Learning}{BADGE Sampling \& Human Review Queue}
    \begin{block}{The Annotation Bottleneck}
        Expert taxonomy time is expensive. Randomly sampling from the unlabelled stream yields redundant, low-diversity annotations.
        \begin{itemize}
            \item \textbf{BADGE Algorithm:} Batch Active learning by Diverse Gradient Embeddings.
            \item \textbf{Dual Criteria:} It computes hypothetical gradient embeddings representing both \textit{epistemic uncertainty} and \textit{feature diversity}.
            \item \textbf{Selection:} Uses $k$-MEANS++ clustering in gradient space to select the optimal batch.
        \end{itemize}
    \end{block}

    \begin{block}{Engineering Rationale}
        By routing only the most mathematically informative samples to Label Studio, we minimise annotation cost while maximising model accuracy improvements per retraining cycle.
    \end{block}

    \vspace{0.2cm}
    \statuspartial

    \script{Annotation by expert taxonomists is expensive and slow. If you send them 10,000 photos of the same common dandelion, you waste their time. BADGE solves this by analyzing the gradients that each unlabelled image would induce in the model. It uses k-MEANS++ clustering in gradient space to pick a batch of images that are both highly uncertain AND visually diverse. This maximizes the model's accuracy improvement per dollar spent on human labeling.}
    \qa{What is the current status of the active learning pipeline?}{The BADGE mathematical sampling algorithm and Label Studio export formats are fully implemented in \texttt{src/taxon\_vision/active\_learning/badge.py}. The automated webhook that triggers unattended model retraining upon annotation completion is our Phase 3 target milestone.}
\end{frame}

\begin{frame}{Active Learning Flow}{Closing the Loop: Stream $\rightarrow$ BADGE $\rightarrow$ Label Studio}
    \begin{center}
    \resizebox{0.95\textwidth}{!}{
    \begin{tikzpicture}[
        node distance=1cm and 1.5cm,
        stepbox/.style={rectangle, draw=deepblue, thick, fill=blue!5, text width=2.5cm, align=center, minimum height=1.2cm, rounded corners},
        arrow/.style={->, thick, >=stealth},
        labelbox/.style={rectangle, draw=signalorange, thick, fill=orange!10, text width=2.5cm, align=center, minimum height=1.2cm, rounded corners}
    ]
        % Nodes
        \node[stepbox] (stream) {\textbf{1. Live Stream} \\ Unlabelled Observations};
        \node[stepbox, right=of stream] (conformal) {\textbf{2. Conformal Filter} \\ Identify low-confidence samples};
        \node[stepbox, right=of conformal] (badge) {\textbf{3. BADGE} \\ Gradient diversity ($k$-MEANS++)};
        \node[labelbox, below=of badge] (labelstudio) {\textbf{4. Label Studio} \\ Expert Human Annotation};
        \node[stepbox, left=of labelstudio] (retrain) {\textbf{5. Retraining Trigger} \\ DVC Pipeline Update};

        % Arrows
        \draw[arrow] (stream) -- (conformal);
        \draw[arrow] (conformal) -- (badge);
        \draw[arrow] (badge) -- (labelstudio);
        \draw[arrow] (labelstudio) -- (retrain);
        \draw[arrow, dashed] (retrain) -| (stream);
    \end{tikzpicture}
    }
    \end{center}

    \script{This flow shows how active learning operates in production. Unlabelled data flows in from our web streams. The split conformal prediction layer filters out cases where the model is highly confident. The remaining low-confidence samples are processed by the BADGE algorithm, which clusters their gradients to select a small, diverse batch. This batch is sent to Label Studio for expert taxonomists to review. Once reviewed, it triggers a DVC pipeline update, retraining the model and updating the live stream.}
    \qa{How does the Conformal Filter help BADGE?}{By filtering out high-confidence samples first, we significantly reduce the size of the dataset BADGE needs to process, making the $O(N \log N)$ $k$-MEANS++ clustering computationally feasible.}
\end{frame}
